\section{Abstract}

This traineeship focused on developing a reproducible finite element breast-model pipeline for the Early Warning Scan concept. The current workflow is implemented as a COMSOL-based Python package in which source-case geometry, material settings, glandular layout, motion input, support assumptions, and lesion parameters are controlled from TOML configuration files. During the final cleanup phase the COMSOL pipeline was made self-contained, so normal COMSOL runs no longer depend on the older FEBio-oriented Python package.

The model was developed in stages. Stage 1 established a mild dynamic motion sanity baseline using a fixed-support acceleration pulse of 0.25g for 0.60 s with mass damping alpha = 60 s$^{-1}$. This case is useful for validating the dynamic input, but it is not a direct anatomical reference for later stages because it uses a larger and simpler baseline geometry. Stage 2 introduced the first fair anatomical baseline: a volume-preserving transverse chestwall curvature with an approximately +55 mm x-offset and automatic nipple/gland alignment. This route produced a breast volume of approximately 585 mL and is used as the current anatomical reference for later stages.

Stage 3 added a more realistic chestwall-aware glandular lobule distribution, increasing the glandular fraction to approximately 24.3\%. Stage 4 then provided a matched no-asymmetry reference for later asymmetry sensitivities, and Stage 5 provided a matched no-Cooper control for later Cooper-ligament support tests. In the current Tier 1 comparison, Stages 3, 4, and 5 are intentionally very similar because the Stage 4 and Stage 5 entries are control cases rather than active sensitivities. The failed default Cooper run means that a reliable Cooper-support effect has not yet been demonstrated.

Stage 6 extends the pipeline toward tumor/lesion sensitivity. The current implementation uses an analytic spherical tumor mask that modifies density and Mooney--Rivlin coefficients locally inside the existing adipose and glandular domains. It does not yet create a separate tumor geometry domain or separate tumor material selection, and should therefore be described as a first-order lesion stiffness/location sensitivity rather than a patient-specific tumor reconstruction.

Overall, the project produced a controlled COMSOL model-development platform rather than a final diagnostic model. The report-ready components are the 0.25g motion definition, the selected Stage 2 chestwall route, the realistic Stage 3 glandular reference as a diagnostic/anatomical baseline, the Stage 4 and Stage 5 control-case structure, and the Stage 6 tumor build-screening workflow. Patient-specific geometry, robust Cooper sensitivity, and quantitative tumor detection remain future work.
